\chapter{결과}
\label{ch:results}

이 장의 수치는 시간순 패치 홀드아웃(학습 15.14 / 검증 15.15 / 시험 15.16)에서 얻은 세 시드 평균의 시험 AUC이며, 논문 \cite{baek2026killconditioned}에 보고된 값이다. 국소화·라벨 수정 후의 공개 코드가 산출하는 수치와의 관계는 제 \ref{sec:disc-audit} 절에서 다룬다. \todo{절제(9.4), SHAP(9.5), 검출기 검증(9.6)의 표는 최종 실행 결과로 채운다.}

\section{기준 성능}
\label{sec:results-baseline}

\begin{table}[htbp]
  \centering
  \caption{입력 표현별 시험 AUC (패치 15.16, 세 시드 평균)}
  \label{tab:baseline}
  \small
  \begin{tabular}{@{}llc@{}}
    \toprule
    모델 & 입력 표현 & 시험 AUC \\
    \midrule
    \textbf{LightGBM} & 공학적 테이블 요약 (약 2{,}980 차원) & \textbf{.675} \\
    MLP & 동일 테이블 요약 (정합 입력) & .626 \\
    BiGRU & 거시 시퀀스 (95 차원 $\times$ 6) & .581 \\
    계층적 융합 & 전역 + 그래프 + 사건 & .581 \\
    Transformer & 거시 시퀀스 & .576 \\
    Cross-Attn & 사건 토큰 $\times$ 플레이어 그래프 & .571 \\
    ST-GNN & 시공간 플레이어 그래프 & .569 \\
    GraphSAGE & 플레이어 그래프 & .569 \\
    \bottomrule
  \end{tabular}
\end{table}

표 \ref{tab:baseline}에서 두 가지 구조적 사실이 드러난다. 첫째, 공학적 테이블 요약 위의 LightGBM이 $.675$로 가장 높고, 정합 입력 MLP가 $.626$으로 그 다음이다. 둘째, 신경망 기반의 시퀀스·그래프·사건 표현은 $.569$--$.581$의 좁은 띠에 모여 있으며, 세 스트림을 모두 결합한 계층적 융합도 단일 BiGRU와 같은 $.581$에 머문다. 즉 표현을 늘리는 것이 신호를 더 드러내지 못했다.

\section{표현 효과와 학습기 효과의 분해}
\label{sec:results-decomp}

RQ2와 RQ3는 표 \ref{tab:baseline}의 두 간격으로 답할 수 있다.

\begin{itemize}[leftmargin=1.5em]
  \item \textbf{표현 효과 (RQ2).} 같은 신경망 학습기 계열 안에서 테이블 요약(MLP, $.626$)과 최선의 비테이블 표현(BiGRU·융합, $.581$)의 차이는 $0.045$이다. 30 초 관측 창이 60 초 프레임 하나 또는 둘만 참조하는 상황에서, 시퀀스·그래프 구조는 계단 유지된 거의 동일한 상태를 여섯 번 반복해서 보는 셈이고, 실제로 변하는 것은 희소한 사건 계수뿐이다. 반면 테이블 요약 $\phi_{\mathrm{tab}}$은 그 희소한 변화를 명시적 통계량(차분, 기울기)으로 노출한다.
  \item \textbf{학습기 효과 (RQ3).} 입력을 고정한 LightGBM($.675$)과 MLP($.626$)의 차이 $0.049$는 순수한 학습기 효과이다. 트리 앙상블은 축 정렬 분할과 잎당 최소 표본 제약으로 수천 차원의 이질적 특징에서 상호작용을 안정적으로 찾는 반면, MLP는 같은 표본 수에서 과적합과 최적화의 어려움을 겪는다. 이는 테이블 데이터에서 그래디언트 부스팅이 심층 신경망보다 우세하다는 광범위한 보고 \cite{grinsztajn2022why}와 일치한다.
\end{itemize}

두 효과의 크기가 비슷하다는 점은 중요하다. 신경망 표현이 뒤처진 이유의 절반은 표현 자체가 아니라 학습기 계열에 있으며, 분 단위 텔레메트리에서는 어떤 표현을 쓰든 \emph{드러낼 수 있는 신호의 총량}이 상한을 이룬다는 해석을 지지한다.

\section{조건부 예측 가능성}
\label{sec:results-conditional}

\begin{table}[htbp]
  \centering
  \caption{하위 집단별 LightGBM 시험 AUC}
  \label{tab:conditional}
  \small
  \begin{tabular}{@{}llc@{}}
    \toprule
    축 & 하위 집단 & AUC \\
    \midrule
    \multirow{3}{*}{게임 국면} & 초반 ($<15$ 분) & .616 \\
     & 중반 (15--25 분) & .712 \\
     & 후반 ($>25$ 분) & .807 \\
    \midrule
    \multirow{3}{*}{온셋 골드 격차} & 대등 ($|\Delta| < 2{,}000$) & .621 \\
     & 보통 & .680 \\
     & 일방 ($|\Delta| > 5{,}000$) & .796 \\
    \bottomrule
  \end{tabular}
\end{table}

표 \ref{tab:conditional}은 RQ4에 답한다. 예측 가능성은 게임 국면에 따라 $.616 \to .712 \to .807$로, 골드 격차에 따라 $.621 \to .680 \to .796$으로 단조 증가한다. 후반과 일방적 교전에서는 누적된 자원 격차가 프레임 수준 상태에 이미 반영되어 있어 분 단위 관측으로도 결과가 거의 결정된다. 반대로 대등한 교전에서는 여러 신경망 표현이 우연 수준($.5$)에 근접한다. 대등한 교전의 결과를 가르는 것은 스킬 적중, 진입 각도, 쿨다운 상태처럼 초 단위 이하에서 일어나는 사건인데, 이는 공개 텔레메트리에 존재하지 않는다. 즉 관측된 AUC 상한은 모델의 한계가 아니라 \emph{정보 입도}의 한계로 해석된다.

\section{도메인 지식 처치의 효과}
\label{sec:results-ablation}

\begin{table}[htbp]
  \centering
  \caption{단일 인자 절제: 처치별 시험 AUC 변화 (\todo{최종 실행 결과로 채울 것})}
  \label{tab:ablation}
  \small
  \begin{tabular}{@{}lcccc@{}}
    \toprule
    처치 & $\Delta_i$ & 95\% CI & DeLong $p$ & Holm 유의 \\
    \midrule
    T1 초점 손실 & -- & -- & -- & -- \\
    T2 게임 국면 부호화 & -- & -- & -- & -- \\
    T3 시간 주의 풀링 & -- & -- & -- & -- \\
    T4 모멘텀 특징 & -- & -- & -- & -- \\
    T5 역할 인식 인접행렬 & -- & -- & -- & -- \\
    T6 다중 과제 손실 & -- & -- & -- & -- \\
    T7 라벨 평활화 & -- & -- & -- & -- \\
    \bottomrule
  \end{tabular}
\end{table}

표 \ref{tab:ablation}은 \texttt{experiment\_runner.py --phase 2 --treatment all}이 산출하는 처치별 한계 효과와 검정 결과를 담는다. 함께 생성되는 산림 그림(forest plot), 상호작용 열지도, 신뢰도 다이어그램을 그림 \ref{fig:forest}에 수록한다. \todo{그림 파일 \texttt{figures/forest.pdf} 추가.}

\begin{figure}[htbp]
  \centering
  \fbox{\parbox{0.8\linewidth}{\centering\vspace{3cm}\todo{처치별 효과 크기 산림 그림}\vspace{3cm}}}
  \caption{처치 T1--T7의 시험 AUC 변화와 95\% 신뢰구간}
  \label{fig:forest}
\end{figure}

\section{특징 기여도 해석}
\label{sec:results-shap}

LightGBM의 TreeSHAP 기여도를 게임 국면별로 층화한 뒤, 특징 이름의 역할(탑·정글·미드·원거리·서포터·팀)과 신호 유형(자원, 스탯, 피해, 사건, 오브젝트)으로 합산하였다. 서포터 역할의 기여 비율은 논문 기준 17.6\%(재실행 시드 7 기준 17.0\%)로, 역할별 비율은 실행에 따라 수 퍼센트포인트 변동한다. \todo{역할 $\times$ 신호 합산 표(\texttt{analysis/shap\_role\_rollup.py} 출력)를 표 \ref{tab:shap}으로 수록.}

\begin{table}[htbp]
  \centering
  \caption{역할별·신호별 SHAP 기여 비율 (\%) (\todo{채울 것})}
  \label{tab:shap}
  \small
  \begin{tabular}{@{}lccccc@{}}
    \toprule
    역할 & 자원 & 스탯 & 피해 & 사건 & 합계 \\
    \midrule
    탑 & -- & -- & -- & -- & -- \\
    정글 & -- & -- & -- & -- & -- \\
    미드 & -- & -- & -- & -- & -- \\
    원거리 & -- & -- & -- & -- & -- \\
    서포터 & -- & -- & -- & -- & -- \\
    팀 전역 & -- & -- & -- & -- & -- \\
    \bottomrule
  \end{tabular}
\end{table}

\section{검출기 검증 결과}
\label{sec:results-detector}

트랙 B(제 \ref{sec:loc-validation} 절)의 기술적 검사는 100 경기에서 480 개 후보 구간(경기당 4.8 개)을 산출하였고 실행 예외, 반복 실행 불일치, 0 킬 구간, 구간 겹침, 미확인 몬스터 사건은 없었다. 사람 주석에 대한 정밀도·재현율·$F_1$·시간 IoU·온셋 오차·킬 없는 교전 누락률은 주석 완료 후 표 \ref{tab:detector}에 보고한다. 같은 프로토콜을 패치 15.14--15.16에도 반복한 뒤에야 시즌 간 강건성을 주장할 수 있다.

\begin{table}[htbp]
  \centering
  \caption{검출기 대 사람 주석 일치도 (\todo{주석 완료 후 채울 것})}
  \label{tab:detector}
  \small
  \begin{tabular}{@{}lcccccc@{}}
    \toprule
    데이터 & 정밀도 & 재현율 & $F_1$ & 평균 IoU & 온셋 오차 (s) & 주석자 간 $\kappa$ \\
    \midrule
    2026 골드셋 (26.13, 100 경기) & -- & -- & -- & -- & -- & -- \\
    2025 (15.14--15.16) & -- & -- & -- & -- & -- & -- \\
    \bottomrule
  \end{tabular}
\end{table}
